A natural end-user question is whether exporting one's listening history and asking a frontier large language model (LLM) for recommendations yields better results than the platform's own recommender. We turn this question into a falsifiable comparison by ranking, for $24$ users with real held-out listening windows, the same $20$-track candidate pool with seven methods: \claude with the user's history, \gpt with the same prompt, \claude without history, popularity, item-kNN, TF-IDF content-based filtering, and a random sanity baseline. We jointly measure rank quality (\recall, \ndcg, \mrr), popularity bias (\arp, \longtail), genre diversity, and free-generation hallucination, following the multi-criterion framework of \citet{epure2025music}. \claude with history attains the highest \ndcg ($0.375$) and ties \gpt while clearly beating \popmethod ($+0.207$ \ndcg, $p{=}0.028$, $d{=}0.53$) and matching \cbf and \itemknn within the resolution of an $N{=}24$ paired Wilcoxon test. Stripping the history collapses \claude below random ($d{=}-1.14$), and adding it back yields a large personalization gain ($+0.29$ \ndcg, $p{=}0.002$, $d{=}0.86$), confirming that the LLM is genuinely conditioning on user data. \claude does not exploit a popularity shortcut: its \arp ($58.6$) is below \cbf's ($59.7$) and well under \popmethod ($74.3$). The result generalizes across LLM families ($\Delta\ndcg < 0.001$ between \claude and \gpt). Overall, off-the-shelf prompting with a Spotify-export-style history matches the strongest classical baseline at a cost of cents and a few seconds per user, but does not significantly exceed it on $N{=}24$ users.
